\chapter*{Prakata}

Buku ini merupakan dokumen yang terus berkembang. Umpan balik, pertanyaan, komentar, dan suntingan dapat dikirimkan ke \href{mailto:open.optimization@gmail.com}{open.optimization@gmail.com}.

Buku ini ditulis dan disusun dengan menggunakan sejumlah materi sumber terbuka. Kami akan berupaya mengutip dengan tepat semua sumber daya yang digunakan serta memberikan rujukan tempat sumber daya tersebut dapat ditemukan. Meskipun materi dalam buku ini berasal dari beragam lisensi, semua materi yang digunakan di sini akan kompatibel dengan CC-BY-SA 4.0; oleh karena itu, seluruh buku ini berada di bawah lisensi CC-BY-SA 4.0.

\begin{center}
\fontsize{14pt}{16pt}\selectfont\textcolor{titletextcolour}{\textbf{UCAPAN TERIMA KASIH UTAMA}}
\end{center}

Saya ingin menyatakan bahwa bagian-bagian penting buku ini diadaptasi dari materi sumber terbuka berikut.
Semua konten yang digunakan kompatibel dengan lisensi CC~BY-SA~4.0 buku ini.

\begin{itemize}
\item \textbf{``A First Course in Linear Algebra''} oleh Lyryx Learning (berdasarkan teks asli karya Ken Kuttler), dilisensikan berdasarkan \textbf{CC~BY~4.0}. Sebagian besar pemformatan mereka digunakan, bersama bagian-bagian terpilih yang membentuk bagian lampiran tentang aljabar linear. Kami sangat berterima kasih kepada Lyryx karena telah membagikan berkas mereka kepada kami. Mereka melakukan pekerjaan luar biasa dalam menyusun buku-buku mereka, dan templat serta pemformatan yang kami gunakan di sini jelas memerlukan banyak kerja untuk disiapkan. Terima kasih telah membagikan seluruh materi ini sehingga penyusunan struktur dan pemformatan buku ini menjadi jauh lebih mudah! Lihat halaman berikutnya untuk daftar kontributor.

\item \textbf{``Foundations of Applied Mathematics''} dengan banyak kontributor, dilisensikan berdasarkan \textbf{CC~BY~3.0~US}. Lihat \url{https://github.com/Foundations-of-Applied-Mathematics}. Beberapa bagian dari catatan ini digunakan beserta sebagian pemformatannya. Sebagian konten telah disunting atau ditata ulang agar sesuai dengan kebutuhan buku ini. Konten ini dilengkapi rujukan kode yang sangat baik dan pemformatan yang bagus untuk menyajikan kode di dalam buku. Lihat halaman berikutnya untuk daftar kontributor.

\item \textbf{``Linear Inequalities and Linear Programming''} oleh
Kevin Cheung, dilisensikan berdasarkan \textbf{CC~BY-SA~4.0}. Lihat \url{https://github.com/dataopt/lineqlpbook}. Catatan ini dipublikasikan di GitHub dalam format ``.Rmd'' agar nyaman dibaca secara daring. Konten ini dikonversi ke \LaTeX\ menggunakan Pandoc. Catatan tersebut membentuk bagian penting dari Bagian Pemrograman Linear buku ini.

\item \textbf{Catatan Pemrograman Linear} oleh Douglas Bish. Catatan ini juga membentuk bagian penting dari Bagian Pemrograman Linear buku ini.

\item \textbf{``Applied Finite Mathematics''} oleh Rupinder Sekhon dan Roberta Bloom, dilisensikan berdasarkan \textbf{CC~BY~4.0}. Lihat \url{https://math.libretexts.org/Bookshelves/Applied_Mathematics/Applied_Finite_Mathematics_(Sekhon_and_Bloom)}. Konten dari Bab~1 (Persamaan Linear) diadaptasi untuk lampiran tentang persamaan linear.
\end{itemize}

Proyek ini didukung oleh Open Education Initiative Faculty Grant dari University Libraries di Virginia Tech. Saya berterima kasih kepada Anita Walz dan Kindred Grey dari University Libraries atas dukungan mereka terhadap proyek ini serta arahan mereka mengenai lisensi terbuka, aksesibilitas, dan penerbitan.
% TODO: ganti/tambahkan pernyataan atribusi hibah yang persis dari
% Anita Walz (wajib dicantumkan dalam bagian awal; lihat surelnya 2026-08).

Saya juga ingin menyampaikan penghargaan kepada Laurent Poirrier dan Diego Moran atas kontribusi berbagai catatan mengenai pemrograman linear dan bilangan bulat.

Saya juga ingin berterima kasih kepada Jamie Fravel yang telah membantu menyunting buku ini serta menyumbangkan bab, contoh, dan kode.

\newpage
